% ╔══════════════════════════════════════════════════════════════╗
% ║              功能筆記 README                                ║
% ║  說明本模板的使用方式、資料夾結構與自訂環境                  ║
% ║  本檔不會被編譯，僅供參考                                    ║
% ╚══════════════════════════════════════════════════════════════╝


% ────────────────────────────────────────────────────────────
%  一、資料夾結構
% ────────────────────────────────────────────────────────────
%
%  專案根目錄/
%  ├── note-template/
%  │   ├── kelong-cjk.sty            ← 字型與數學字族設定（請勿隨意修改）
%  │   ├── fonts/                    ← 字型檔，必須與 main.tex 同層
%  │   ├── main.tex                  ← 編譯入口，切換文件只改這裡
%  │   ├── package_and_settings.tex  ← 所有套件與全域設定
%  │   └── 功能筆記readme.tex        ← 本說明文件
%  │
%  └── content/                      ← 所有文件內容
%      └── {科目}/                   ← 例如 linear_algebra
%          ├── note/note_1/          ← 筆記，依章節編號
%          ├── hw/hw_1/              ← 作業，依次序編號
%          ├── quiz/quiz_1/          ← 小考
%          ├── exam/                 ← 考試
%          │   ├── midterm/          ← 期中考
%          │   └── final/            ← 期末考
%          └── extra_hw/extra_hw_1/  ← 補充習題
%
%  ► 新增科目：在 content/ 下建立新資料夾，結構與上方相同。
%  ► 新增章節筆記：複製 note/note_1/ 整個資料夾，改名為 note_2/，
%                  修改 part0_all.tex 的 \input 路徑即可。


% ────────────────────────────────────────────────────────────
%  二、每個文件資料夾的內部結構
% ────────────────────────────────────────────────────────────
%
%  {文件資料夾}/
%  ├── 0. 標題與作者.tex   ← 定義 \mytitle / \myauthor / \title / \author
%  ├── part0_all.tex        ← 依序 \input 所有 part 檔的索引
%  ├── part1.tex            ← 第一段內容
%  ├── part2.tex            ← 第二段內容（視需要新增）
%  └── figures/             ← 本文件用到的圖片


% ────────────────────────────────────────────────────────────
%  三、圖片管理規則
% ────────────────────────────────────────────────────────────
%
%  圖片放在各文件的 figures/ 子資料夾：
%    content/{科目}/note/note_1/figures/my_diagram.png
%
%  引用時使用從根目錄出發的完整相對路徑（不要只寫檔名）：
%
%    \begin{figure}[H]
%      \centering
%      \includegraphics[width=0.8\textwidth]
%        {content/{科目}/note/note_1/figures/my_diagram}
%      \caption{圖片標題}
%    \end{figure}
%
%  ► 注意：本模板沒有設全域 \graphicspath，必須使用完整路徑，
%           否則 Overleaf 找不到圖片。
%  ► 命名建議：使用說明性英文名稱（如 ols_scatter.png），
%               避免 image1.png 等無意義名稱。


% ────────────────────────────────────────────────────────────
%  四、如何切換文件（main.tex 的操作）
% ────────────────────────────────────────────────────────────
%
%  只需修改 main.tex 中的兩行 \input：
%
%    % 步驟一：標題與作者（文件開頭，\begin{document} 之前）
%    \input{content/{科目}/{類型}/{文件名}/0. 標題與作者}
%
%    % 步驟二：文件內容（在 \begin{document} 內）
%    \input{content/{科目}/{類型}/{文件名}/part0_all}
%
%  範例：
%    \input{content/linear_algebra/note/note_1/0. 標題與作者}
%    \input{content/linear_algebra/note/note_1/part0_all}


% ────────────────────────────────────────────────────────────
%  五、自訂環境使用說明
% ────────────────────────────────────────────────────────────

% ── narrowframed（定義、定理等重要內容）────────────────────
%
%  \begin{narrowframed}
%    \textbf{定義 1.1.1} 定義名稱\\
%    定義內容...
%  \end{narrowframed}

% ── whitebox（無框視覺隔離區塊）────────────────────────────
%
%  \begin{whitebox}
%    內容...
%  \end{whitebox}

% ── 習題三件套（題目 → 解答 → 分析）────────────────────────
%
%  標準順序：examplebox → solutionbox → analysisbox（可省略）
%
%  \begin{examplebox}{題號}{題目名稱}
%    題目敘述...
%  \end{examplebox}
%
%  \begin{solutionbox}
%    解題過程...
%  \end{solutionbox}
%
%  \begin{analysisbox}
%    解題思路、考點說明...（接在 solutionbox 之後）
%  \end{analysisbox}
%
%  ► analysisbox 與 narrowframed 外觀相同（淡藍底 90% 寬方框），
%    框內頂端自動顯示「題目分析」標題，字體大小與「解答」相同。

% ── 程式碼區塊 ─────────────────────────────────────────────
%
%  \begin{lstlisting}[language=R]
%    lm_model <- lm(y ~ x, data = df)
%    summary(lm_model)
%  \end{lstlisting}
%
%  行內程式碼：\lstinline|your_code|

% ── 多欄排版 ───────────────────────────────────────────────
%
%  \begin{multicols}{2}
%    左欄內容...
%    \columnbreak
%    右欄內容...
%  \end{multicols}

% ── 字型指令 ───────────────────────────────────────────────
%
%  \kai{楷體文字}    \yuan{圓體文字}    \hei{黑體文字}
%  （正文預設為明體，粗體 \textbf{} 自動切換為黑體）
